\clearpage
\VerdierPhysicalPageAnchor{0172}%
\begingroup\small\setlength{\parskip}{0pt}\setlength{\emergencystretch}{1.5em}
\VerdierSemanticAnchor{ch4-4-4}
\medskip\noindent\textbf{4.4. 収束。}

\VerdierSemanticAnchor{ch4-4-4-1}
\medskip\noindent\textbf{4.4.1.}\quad
スペクトル系列の収束に関する詳しい研究については\emph{同所引用文献}を
参照されたい。ここでは次の判定条件だけを用いる：

\VerdierSemanticAnchor{ch4-4-4-2}
\medskip\noindent\textbf{定義 4.4.2.}\quad
スペクトル対象 $(H^n(p,q),\delta)$ がアーベル圏 $\mathcal A$ に値をとる
とする。この対象が\emph{定常}であるとは、各整数 $n$ に対して二つの整数
$u(n),v(n)\in\mathbb Z$ が存在して、次を満たすことをいう：
\[
\begin{aligned}
H^n(p,q)&=0,&&u(n)\leq p\leq q\leq+\infty,\\
H^n(p,q)&=0,&&-\infty\leq p\leq q\leq v(n).
\end{aligned}
\]

\VerdierSemanticAnchor{ch4-4-4-3}
\medskip\noindent\textbf{定理 4.4.3.}\quad
\textit{$(H^n(p,q),\delta)$ を定常スペクトル対象とする。}

\textit{a) 各 $n$ に対して、フィルトレーション
$F^pH^n(-\infty,+\infty)$、$p\in\widetilde{\mathbb Z}$ は定常である。
すなわち、二つの整数 $n_1\leq n_2$ が $\mathbb Z$ に存在して、次を満たす：}
\[
F^{n_1}H^n(-\infty,+\infty)=0
\quad\text{かつ}\quad
F^{n_2}H^n(-\infty,+\infty)=H^n(-\infty,+\infty).
\]

\textit{b)
$E\!\left[\begin{smallmatrix}n\\r,p,q,s\end{smallmatrix}\right]$ から出る
（それぞれ、そこに入る）微分は、$r\leq v(n+1)$（それぞれ、
$u(n-1)\leq s$）ならば零であり、さらに次が成り立つ：}

\textit{c) 各 $n$ および各組 $p\leq q$ に対して、対象
$E\!\left[\begin{smallmatrix}n\\r,p,q,s\end{smallmatrix}\right]$ のうち、}
\[
-\infty\leq r\leq v(n+1)
\quad\text{かつ}\quad
u(n-1)\leq s\leq+\infty
\]
\textit{を満たすものは標準同型である。}

この定理は定義から直ちに従う。

\VerdierSemanticAnchor{ch4-4-4-3-convergence}
定常スペクトル対象は、\emph{同所引用文献}の意味で強収束するスペクトル系列を
与える。スペクトル対象に付随するスペクトル系列について、そのスペクトル対象
が定常であるとき、その系列を\emph{定常}、または語を広く用いて\emph{収束する}
という。(4.3.3) のいずれかの記法を用いて、次のように書く：
\[
\mathrm I_r^{p,q}\Longrightarrow H^{p+q}(-\infty,+\infty),
\qquad
\mathrm{II}_r^{p,q}\Longrightarrow H^{p+q}(-\infty,+\infty)
\]
これはスペクトル系列が収束することを表す。このとき、項
$H^{p+q}(-\infty,+\infty)$ の全体をスペクトル系列の\emph{収束先}と呼ぶ。
\endgroup
